\documentclass[../../../include/open-logic-section]{subfiles}

\begin{document}

\olfileid{sth}{cardinals}{hp}
\olsection{પરિશિષ્ટ: હ્યુમનો સિદ્ધાંત}

\olref[cp]{sec}માં આપણે કૅન્ટરનો સિદ્ધાંત વર્ણવ્યો હતો.
તે આ હતો:
\begin{align*}
	\card{A} = \card{B} & \text{ iff } A \approx B.
\intertext{આ આજકાલ \emph{હ્યુમનો સિદ્ધાંત} કહેવાય છે
તેના જેવો છે; તે કહે છે:}
	\fregenum{x} {F(x)} = \fregenum{x}{G(x)} & \text{ iff } F \sim G
\end{align*}
અહીં `$F \sim G$' એટલે $F$ હોય એવી વસ્તુઓ
$G$ હોય એવી વસ્તુઓ જેટલી જ છે; એટલે $F$ હોય
એવી વસ્તુઓ અને $G$ હોય એવી વસ્તુઓ વચ્ચે એક-એક
અને વ્યાપ્ત વિધેય મૂકી શકાય છે; ઔપચારિક રીતે:
\begin{align*}
	\exists R(&\forall v\forall y(Rvy \lif (Fv \land Gy)) \land {}\\
		&\forall v(Fv \lif \lexists![y][Rvy]) \land {}\\
		&\forall y(Gy \lif \lexists![v][Rvy]))
\end{align*}
પણ હ્યુમના સિદ્ધાંત અને કૅન્ટરના સિદ્ધાંત વચ્ચે
પ્રકારનો ભેદ છે. કૅન્ટરના સિદ્ધાંતના વિધાનમાં
``$A$'' અને ``$B$'' પ્રથમ-ક્રમના પદો છે અને
\emph{ગણો} દર્શાવે છે. હ્યુમના સિદ્ધાંતના વિધાનમાં
``$F$'', ``$G$'' અને ``$R$'' \emph{પ્રથમ-ક્રમના પદો
નથી}; તેઓ \emph{વિધેયકના સ્થાને} છે. (કદાચ તેઓ
\emph{ગુણધર્મો} દર્શાવે છે.) તેથી હ્યુમના સિદ્ધાંતનો
અર્થ આમ કહી શકાય: $F$ હોય એવી વસ્તુઓની સંખ્યા
$G$ હોય એવી વસ્તુઓની સંખ્યા જેટલી જ હોય જો અને
માત્ર જો $F$ હોય એવી વસ્તુઓને $G$ હોય એવી
વસ્તુઓ સાથે એક-એક રીતે જોડી શકાય. તેને
\emph{હ્યુમનો સિદ્ધાંત} કહે છે, કારણ કે હ્યુમે
એક વાર આ લખ્યું હતું:
\begin{quote}
  બે સંખ્યાઓને એવી રીતે જોડવામાં આવે કે એકની દરેક
  એકમને બીજીની એક એકમ અનુરૂપ હોય, તો આપણે તેમને
  સમાન કહીએ છીએ.
  \citep[Pt.III Bk.1 \S1]{Hume1740}
\end{quote}
આ અવતરણ ટાંકીને પછી (અર્થની દૃષ્ટિએ) હ્યુમના
સિદ્ધાંતની રૂપરેખા આપનાર \citet[\S63]{Frege1884}એ
તેને આધુનિક ગણિતીય-તાર્કિક ચર્ચામાં અગત્યનું સ્થાન
આપ્યું હતું.

હ્યુમના સિદ્ધાંત અને મૂળભૂત નિયમ~Vની રચનાત્મક
સમાનતા નોંધો. તેને આપણે \olref[story][blv]{sec}માં
આમ રજૂ કર્યો હતો:
\[
	\fregeext{x}{F(x)} = \fregeext{x}{G(x)} \text{iff } \lforall[x][(F(x) \liff G(x))].
\]
અહીં થોડી ટિપ્પણી અને સરખામણી ઉપયોગી થશે.

હ્યુમના સિદ્ધાંત કે મૂળભૂત નિયમ~V જેવા સિદ્ધાંતને
લેવાની બે રીતો છે: \emph{વિધેયક} અથવા
\emph{અવિધેયક} (\olref[story][predicative]{sec}
યાદ કરો). મૂળભૂત નિયમ~Vને અવિધેયક રીતે લઈએ,
તો દરેક~$F$ માટે $\fregeext{x}{F(x)}$ વસ્તુ
એ જ પરિમાણન-ક્ષેત્રમાં આવે છે જેનો નિયમ~V
ઘડવા વાપર્યો હતો. એ જ રીતે હ્યુમના સિદ્ધાંતને
અવિધેયક રીતે લઈએ, તો દરેક~$F$ માટે
$\fregenum{x}{F(x)}$ વસ્તુ એ જ પરિમાણન-ક્ષેત્રમાં
આવે છે જેનો હ્યુમનો સિદ્ધાંત ઘડવા વાપર્યો હતો.
તેનાથી વિપરીત, \emph{વિધેયક} અર્થમાં
$\fregeext{x}{F(x)}$ અને~$\fregenum{x}{F(x)}$
વસ્તુઓ કોઈ \emph{જુદા} ક્ષેત્રની હશે.

હવે મૂળભૂત નિયમ~Vને અવિધેયક રીતે વાંચીએ, તો
સહજ ગણરચના દ્વારા અસુસંગતતા મળે છે
(વિગતો માટે \olref[story][blv]{sec} જુઓ).
સહજ ગણરચનાની જેમ તેને \emph{વિધેયક} રીતે
વાંચવાથી સુસંગત બનાવી શકાય છે. પરંતુ ત્યારે તે
કદાચ આપણે ઇચ્છેલી બધી બાબતો પૂરી નહીં પાડે.

હ્યુમનો સિદ્ધાંત, જોકે, સુસંગત રીતે
\emph{અવિધેયક} વાંચી શકાય છે. એમ વાંચતાં તે
ઘણો શક્તિશાળી બને છે.

દાખલા તરીકે ``$x \neq x$'' વિધેયક લો, જેને
દેખીતી રીતે કોઈ વસ્તુ સંતોષતી નથી. હ્યુમના
સિદ્ધાંતથી હવે $\# x( x\neq x)$ વસ્તુ મળે છે.
તેને સંખ્યા~$0$ લઈ શકીએ. હવે \emph{અવિધેયક}
અર્થમાં---અને \emph{ફક્ત} એ અર્થમાં---આ વસ્તુ
$0$ મૂળ પરિમાણન-ક્ષેત્રમાં આવે છે. તેથી
``$x = 0$'' વિધેયક સાથે હ્યુમનો સિદ્ધાંત લગાડી
$\#x (x = 0)$ વસ્તુ મેળવી શકીએ. તેને સંખ્યા~$1$
લઈ શકીએ. વળી હ્યુમના સિદ્ધાંતથી $0 \neq 1$
મળે છે, કારણ કે પોતાનાથી અસમાન વસ્તુઓ અને $0$
સાથે સમાન વસ્તુઓ વચ્ચે એક-એક અને વ્યાપ્ત વિધેય
હોઈ શકે નહીં (પહેલા પ્રકારની કોઈ વસ્તુ નથી, પણ
બીજા પ્રકારની એક છે). ફરી અવિધેયક રીતે કામ કરતાં
$1$ પણ મૂળ પરિમાણન-ક્ષેત્રમાં આવે છે. તેથી
``$(x = 0 \lor x = 1)$'' વિધેયક સાથે હ્યુમનો
સિદ્ધાંત લગાડી $\#x(x = 0 \lor x = 1)$ વસ્તુ
મેળવી શકીએ. તેને સંખ્યા~$2$ લઈ શકીએ અને
$0\neq 2$ તથા $1 \neq 2$ વગેરે બતાવી શકીએ.

આમ, અવિધેયક રીતે લેવાયેલ હ્યુમના સિદ્ધાંતથી
\emph{અનંત ઘણી વસ્તુઓ} અસ્તિત્વમાં છે તે મળે છે.
આ કારણે \emph{નવા ફ્રેગેવાદી તર્કવાદીઓ}એ
હ્યુમના સિદ્ધાંતને અંકગણિતનો પાયો લેવાનું પસંદ કર્યું છે.

જોકે ફ્રેગેએ \emph{પોતે} હ્યુમના સિદ્ધાંતને
અંકગણિતનો પાયો લીધો નહોતો. તેના બદલે ફ્રેગેએ
એક સ્પષ્ટ વ્યાખ્યામાંથી હ્યુમનો સિદ્ધાંત સાબિત
કર્યો: $\fregenum{x}{F(x)}$ને $F \sim \Phi$
ખ્યાલના વિસ્તાર તરીકે વ્યાખ્યાયિત કર્યો.
આધુનિક સંકેતમાં તેને
$\fregenum{x}{F(x)} = \Setabs{G}{F \sim G}$
રૂપે દર્શાવવાનો પ્રયાસ કરી શકીએ; પરંતુ એ આપણને
ફરી સહજ ગણરચનાની મુશ્કેલીઓમાં લઈ જશે.

\end{document}
